% BTS CIEL 1 — Réseaux — 06 — DHCP
% Compiler deux fois avec XeLaTeX depuis fichiers/.
\documentclass[aspectratio=169,11pt]{beamer}
\usepackage{fontspec}
\setmainfont{Latin Modern Sans}
\setsansfont{Latin Modern Sans}
\setmonofont{Latin Modern Mono}
\usepackage{polyglossia}
\setmainlanguage{french}
\usepackage{tikz}
\usetikzlibrary{arrows.meta,positioning,calc,shapes.geometric}
\usepackage{listings}
\usepackage{xcolor}
\definecolor{foret}{HTML}{1C322D}
\definecolor{creme}{HTML}{F8F3EE}
\definecolor{or}{HTML}{EBB552}
\definecolor{rose}{HTML}{F1CDBE}
\definecolor{vertclair}{HTML}{DDE9DF}
\definecolor{encre}{HTML}{24372E}
\setbeamercolor{normal text}{fg=creme,bg=foret}
\setbeamercolor{frametitle}{fg=or,bg=foret}
\setbeamercolor{title}{fg=creme}
\setbeamercolor{block title}{fg=foret,bg=or}
\setbeamercolor{block body}{fg=encre,bg=creme}
\setbeamercolor{alerted text}{fg=or}
\setbeamercolor{structure}{fg=or}
\setbeamertemplate{navigation symbols}{}
\setbeamertemplate{footline}{\hfill\color{or}\insertframenumber\hspace{0.5cm}\vspace{0.25cm}}
\setbeamertemplate{frametitle}{\vspace{0.35cm}\insertframetitle\par\vspace{0.15cm}\color{or}\rule{\linewidth}{0.6pt}}
\setlength{\parskip}{5pt}
\lstdefinestyle{webcode}{basicstyle=\ttfamily\small\color{encre},keywordstyle=\color{foret},stringstyle=\color{foret},
  backgroundcolor=\color{creme},frame=single,rulecolor=\color{or},framerule=1pt,
  xleftmargin=5pt,xrightmargin=5pt,columns=fullflexible,keepspaces=true,showstringspaces=false}
\newcommand{\question}[1]{\begin{beamercolorbox}[rounded=true,sep=12pt]{block body}\Large\textbf{#1}\end{beamercolorbox}}
\newcommand{\etiquette}[1]{{\color{or}\small\MakeUppercase{#1}}}
% Les schémas ont une surface claire, également lisible en noir et blanc.
\tikzset{
  boite/.style={draw=encre,fill=creme,text=encre,rounded corners,
    align=center,text width=2.8cm,minimum height=1cm},
  flux/.style={-{Stealth},thick,draw=encre},
  every node/.style={text=encre,font=\small}}
\newenvironment{schema}{%
  \begin{beamercolorbox}[sep=8pt,rounded=true]{block body}
  \centering\begin{tikzpicture}[>=Stealth]
}{\end{tikzpicture}\end{beamercolorbox}}
\title{Le protocole DHCP}
\subtitle{BTS CIEL 1 · Réseaux · 06 — DHCP}
\author{} \date{}
\begin{document}
\begin{frame}[plain]
\etiquette{Configurer automatiquement les postes}\par\bigskip
{\Huge\bfseries Le protocole DHCP}\par\bigskip
{\Large BTS CIEL 1 · Réseaux · 06 — DHCP}\par\bigskip
\begin{schema}
\node[boite,text width=4cm,minimum height=2cm,font=\large] (a) {Poste};
\node[boite,text width=4cm,minimum height=2cm,font=\large,right=4cm of a] (b) {Serveur DHCP};
\draw[flux] ([yshift=5mm]a.east)--node[above]{demande} ([yshift=5mm]b.west);
\draw[flux] ([yshift=-5mm]b.west)--node[below]{configuration} ([yshift=-5mm]a.east);
\end{schema}
\end{frame}
\begin{frame}{Trente postes à configurer}
\question{Comment éviter les erreurs de saisie et les adresses en double ?}
\begin{schema}
\node[boite,text width=3cm,minimum height=2.5cm,font=\normalsize] (a) at (0,0) {30 postes\\4 paramètres chacun};
\node[boite,text width=3cm,minimum height=2.5cm,font=\normalsize] (b) at (5,0) {Saisie manuelle\\répétée};
\node[boite,text width=3cm,minimum height=2.5cm,font=\normalsize] (c) at (10,0) {Erreurs ou\\conflits d'adresses};
\draw[flux] (a)--(b); \draw[flux] (b)--(c);
\end{schema}
\end{frame}
\begin{frame}{Une adresse prêtée comme une chambre}
\question{Comment attribuer une chambre disponible pour un séjour ?}
\begin{schema}
\node[boite,text width=4cm,minimum height=2.5cm,font=\normalsize] (a) at (0,0) {Client de l'hôtel\\demande une chambre};
\node[boite,text width=4cm,minimum height=2.5cm,font=\normalsize] (b) at (8,0) {Accueil\\prête une clé\\pour une durée};
\draw[flux] ([yshift=5mm]a.east)--node[above]{1. demande} ([yshift=5mm]b.west);
\draw[flux,dashed] ([yshift=-5mm]b.west)--node[below]{2. attribution} ([yshift=-5mm]a.east);
\end{schema}
DHCP prête une adresse pour une durée : le bail.
\end{frame}
\begin{frame}{DORA : un poste obtient son adresse}
Réseau \texttt{192.168.10.0/24} ; serveur \texttt{192.168.10.2}.
\begin{schema}
\node[font=\normalsize] at (0,3.3) {Poste sans adresse};
\node[font=\normalsize] at (9,3.3) {Serveur DHCP};
\draw[dashed] (0,3.0)--(0,.1); \draw[dashed] (9,3.0)--(9,.1);
\draw[flux,dashed] (0,2.7)--node[above,fill=creme,inner sep=2pt]{1. Discover : recherche diffusée} (9,2.7);
\draw[flux] (9,1.9)--node[above,fill=creme,inner sep=2pt]{2. Offer : propose 192.168.10.100} (0,1.9);
\draw[flux,dashed] (0,1.1)--node[above,fill=creme,inner sep=2pt]{3. Request : choix diffusé} (9,1.1);
\draw[flux] (9,.3)--node[above,fill=creme,inner sep=2pt]{4. Ack : confirme le bail} (0,.3);
\end{schema}
{\small Pointillés : diffusion, à tous les postes du réseau local ;\\
trait plein : envoi au seul poste.\\
Offer et Ack : diffusés ou envoyés au seul poste, selon le client.\\
DHCP utilise UDP : client port 68, serveur port 67.}
\end{frame}
\begin{frame}{Quatre paramètres pour communiquer}
\begin{beamercolorbox}[sep=12pt,rounded=true]{block body}
\centering\renewcommand{\arraystretch}{1.7}
\begin{tabular}{lll}
\textbf{Paramètre} & \textbf{Exemple} & \textbf{Rôle}\\\hline
Adresse IPv4 & \texttt{192.168.10.100} & Identifier le poste\\
Masque & \texttt{255.255.255.0} & Repérer le réseau local\\
Passerelle & \texttt{192.168.10.1} & Joindre d'autres réseaux\\
Serveur DNS & \texttt{192.168.10.2} & Résoudre les noms\\
\end{tabular}
\end{beamercolorbox}
DHCP fournit l'adresse du serveur DNS ; DNS répond aux demandes de noms.
\end{frame}
\begin{frame}{Réserver les adresses fixes}
\begin{schema}
\draw[thick] (0,0)--(11,0);
\foreach \x/\label in {0/.1,1/.2,5/.100,8/.119,11/.254}
  \draw (\x,-.12)--(\x,.12) node[below=4pt] {\label};
\draw[line width=3pt] (0,.35)--(1,.35);
\node[above,align=center] at (.5,.45) {adresses fixes\\(.1, .2)};
\draw[line width=3pt] (5,.35)--(8,.35);
\node[above] at (6.5,.45) {pool .100 à .119};
\end{schema}
La passerelle utilise \texttt{192.168.10.1}.\\
Le serveur DHCP et DNS utilise \texttt{192.168.10.2} : il assure les deux rôles.\\
Le pool est la plage à prêter. Les adresses fixes restent hors du pool.
\end{frame}
\begin{frame}{Compter les adresses du pool}
\begin{schema}
\node[boite,text width=3cm,minimum height=2cm,font=\normalsize] (a) at (0,0) {Début : .100};
\node[boite,text width=3cm,minimum height=2cm,font=\normalsize] (b) at (5,0) {Fin : .119};
\node[boite,text width=3cm,minimum height=2cm,font=\normalsize] (c) at (10,0) {20 adresses};
\draw[flux] (a)--(b); \draw[flux] (b)--(c);
\end{schema}
\begin{block}{Méthode : plage dans un même /24}
1. Relevez le dernier nombre de chaque adresse.\\
2. Calculez fin $-$ début.\\
3. Ajoutez 1 pour compter les deux extrémités.\\
Exemple : $119-100+1=20$.
\end{block}
\end{frame}
\begin{frame}{Le bail se renouvelle}
\begin{schema}
\node[boite,text width=2.5cm,minimum height=2.5cm,font=\small] (a) at (0,0) {Adresse attribuée\\pour 8 heures};
\node[boite,text width=2.5cm,minimum height=2.5cm,font=\small] (b) at (5,0) {Le poste demande\\à renouveler};
\node[boite,text width=2.5cm,minimum height=2.5cm,font=\small] (c) at (10,0) {Le serveur confirme\\un nouveau bail};
\draw[flux] (a)--node[above,align=center,font=\scriptsize]{1. à mi-bail\\(4 h)} (b); \draw[flux] (b)--node[above,align=center,font=\scriptsize]{2. accord :\\8 h de plus} (c);
\end{schema}
Sans renouvellement avant expiration, le poste cesse d'utiliser l'adresse.
\end{frame}
\begin{frame}{Renouveler sous Windows}
\begin{schema}
\node[boite,text width=4cm,minimum height=2.5cm,font=\normalsize] (a) at (0,0) {1. Libérer le bail\\\texttt{ipconfig /release}};
\node[boite,text width=4cm,minimum height=2.5cm,font=\normalsize] (b) at (8,0) {2. Demander à nouveau\\\texttt{ipconfig /renew}};
\draw[flux] (a)--(b);
\end{schema}
Dans l'invite de commandes Windows, ces commandes agissent sur les cartes réseau configurées par DHCP.
\end{frame}
\begin{frame}{Pas de réponse : vérifier la configuration}
\begin{schema}
\node[boite,text width=3cm] (a) at (0,0) {Poste\\169.254.12.7};
\node[boite,text width=3cm] (b) at (8,0) {Serveur DHCP\\service arrêté};
\draw[flux,dashed] (a)--node[above]{Discover sans réponse} (b);
\end{schema}
{\small Une adresse APIPA est une adresse que Windows se donne seul, faute de réponse DHCP. Le poste ne joint que ses voisins en 169.254 ; sans passerelle, il ne sort pas du réseau local.}
\begin{block}{Méthode}
1. Repérez une adresse en \texttt{169.254}.\\
2. Vérifiez le câble et les voyants.\\
3. Vérifiez que le service DHCP du serveur est actif.\\
\quad Packet Tracer : serveur > Services > DHCP > On.\\
4. Relancez \texttt{ipconfig /renew} dans Windows.
\end{block}
\end{frame}
\begin{frame}{À vous : interpréter un échange}
\question{Le poste peut-il utiliser l'adresse proposée ? Pourquoi ?}
\begin{schema}
\node[boite,text width=3cm,minimum height=2.5cm,font=\normalsize] (a) at (0,0) {Poste\\Discover envoyé};
\node[boite,text width=4cm,minimum height=2.5cm,font=\normalsize] (b) at (8,0) {Serveur\\Offer : \texttt{192.168.10.105}};
\draw[flux,dashed] ([yshift=5mm]a.east)--node[above]{1. Discover} ([yshift=5mm]b.west);
\draw[flux] ([yshift=-5mm]b.west)--node[below]{2. Offer} ([yshift=-5mm]a.east);
\end{schema}
\end{frame}
\begin{frame}{Correction : attendre la confirmation}
\begin{schema}
\node[boite,text width=3cm,minimum height=2.5cm,font=\normalsize] (a) at (0,0) {Offer\\proposition};
\node[boite,text width=3cm,minimum height=2.5cm,font=\normalsize] (b) at (5,0) {Request\\choix du client};
\node[boite,text width=3cm,minimum height=2.5cm,font=\normalsize] (c) at (10,0) {Ack\\confirmation};
\draw[flux] (a)--(b); \draw[flux] (b)--(c);
\end{schema}
Plusieurs serveurs peuvent faire une offre. Le Request annonce le choix du poste ; l’Ack confirme le bail.
\end{frame}
\begin{frame}{À vous : dimensionner une plage}
\question{Combien d'adresses cette plage contient-elle ?}
\begin{schema}
\node[boite,text width=4cm,minimum height=2.5cm,font=\normalsize] (a) at (0,0) {Première adresse\\\texttt{192.168.20.40}};
\node[boite,text width=4cm,minimum height=2.5cm,font=\normalsize] (b) at (8,0) {Dernière adresse\\\texttt{192.168.20.69}};
\draw[flux] (a)--node[above]{plage inclusive} (b);
\end{schema}
\end{frame}
\begin{frame}{Correction : compter les deux extrémités}
\begin{schema}
\node[boite,text width=3cm,minimum height=2.5cm,font=\normalsize] (a) at (0,0) {Fin : 69\\Début : 40};
\node[boite,text width=3cm,minimum height=2.5cm,font=\normalsize] (b) at (5,0) {$69-40=29$\\écarts};
\node[boite,text width=3cm,minimum height=2.5cm,font=\normalsize] (c) at (10,0) {$29+1=30$\\adresses};
\draw[flux] (a)--(b); \draw[flux] (b)--(c);
\end{schema}
Les adresses de début et de fin sont toutes deux comptées.
\end{frame}
\begin{frame}{À retenir}
\begin{schema}
\node[boite,text width=3cm,minimum height=1.8cm,font=\normalsize] (a) at (0,0) {Poste};
\node[boite,text width=4cm,minimum height=1.8cm,font=\normalsize] (b) at (9,0) {Serveur DHCP\\pool et baux};
\draw[flux,<->] (a)--node[above,font=\small]{DORA : attribution initiale} (b);
\end{schema}
\begin{block}{Cinq repères}
DHCP fournit les paramètres IPv4.\\
DORA : Discover, Offer, Request, Ack.\\
DHCP utilise UDP : client port 68, serveur port 67.\\
Le pool exclut les adresses fixes ; le bail a une durée.\\
Une adresse APIPA peut signaler l'absence de réponse DHCP.
\end{block}
\end{frame}
\begin{frame}{Sources}
\begin{beamercolorbox}[sep=12pt,rounded=true]{block body}
\textbf{RFC 2131} — attribution et bail DHCP\\
\texttt{https://www.rfc-editor.org/rfc/rfc2131}\\[6pt]
\textbf{RFC 2132} — options de configuration\\
\texttt{https://www.rfc-editor.org/rfc/rfc2132}\\[6pt]
\textbf{RFC 3927} — adresses locales IPv4\\
\texttt{https://www.rfc-editor.org/rfc/rfc3927}\\[6pt]
\textbf{Microsoft Learn} — commande ipconfig\\
{\scriptsize\texttt{https://learn.microsoft.com/windows-server/administration/windows-commands/ipconfig}}\\[6pt]
Support enseignant : « TP DHCP, HTTP, DNS, FTP », p. 1--6.
\end{beamercolorbox}
\end{frame}
\end{document}
